\chapter{Semantic Types, Tags, and Typed Morphisms}
\label{ch:semantic-types}

This chapter formalizes the \emph{typing system} underlying Polyxan hypergraphs, together with the morphisms and rewrite rules that preserve semantic consistency. Building on Chapter~\ref{ch:polyxan-atoms}, we introduce semantic types, tag refinements, and typed morphisms; these form the foundation for type-aware rewrites, constraint propagation, and Polycompiler correctness (Chapter~\ref{ch:polycompiler}).

% ============================================================
\section{Semantic Types and Type Assignments}
% ============================================================

Let $\mathcal{G} = (V,E)$ be a Polyxan hypergraph (Chapter~\ref{ch:polyxan-atoms}), where $V$ is the set of content atoms and $E$ is the set of hyperedges encoding relations.

\begin{definition}[Semantic Type]
A \emph{semantic type} is an element of a type category $\mathbf{Type}$, assigned to all vertices and hyperedges:
\[
\mathrm{type} : V \cup E \to \mathbf{Type}.
\]
Types encode compositional constraints, inheritance, semantic roles, and permitted relational patterns.
\end{definition}

\begin{example}
Vertices might have types
\[
\texttt{Matter},\quad \texttt{Energy},\quad \texttt{Field},
\]
while hyperedges could have types such as
\[
\texttt{causes},\quad \texttt{interactsWith},\quad \texttt{encodes}.
\]
\end{example}

These assignments enforce both structural and semantic validity across the Polyxan hypergraph.

% ============================================================
\section{Tags as Type Refinements}
% ============================================================

\begin{definition}[Semantic Tag]
A \emph{tag} is a refinement of a semantic type, drawn from a tag set $\mathbf{Tag}$ equipped with a projection to the base type:
\[
\mathrm{tag}: V \cup E \to \mathbf{Tag},
\qquad
\pi_\tau : \mathbf{Tag} \to \mathbf{Type}.
\]
Tags encode subtyping information, contextual attributes, epistemic metadata, or runtime embeddings.
\end{definition}

\begin{remark}
Tags often record simulation-relevant metadata:
confidence scores, domain embeddings, temporal qualifiers, symbolic lineage, or cross-modal annotations.
\end{remark}

Type compatibility requires:
\[
\pi_\tau(\mathrm{tag}(x)) = \mathrm{type}(x),
\]
ensuring every tag refines a compatible semantic type.

% ============================================================
\section{Typed Morphisms}
% ============================================================

A structure-preserving mapping between typed Polyxan hypergraphs must also maintain semantic types and tags.

\begin{definition}[Typed Hypergraph Morphism]
Let $\mathcal{G}=(V,E)$ and $\mathcal{G}'=(V',E')$ be typed Polyxan hypergraphs with type/tag assignments $\mathrm{type},\mathrm{tag}$ and $\mathrm{type}',\mathrm{tag}'$.  
A \emph{typed morphism} is a pair of functions:
\[
f_V : V \to V',
\qquad
f_E : E \to E',
\]
such that:

\begin{enumerate}
\item \textbf{Structure preservation:}  
For each hyperedge $e = \{v_1,\dots,v_k\} \in E$,
\[
f_E(e) = \{f_V(v_1),\dots,f_V(v_k)\} \in E'.
\]

\item \textbf{Type refinement:}  
For all $x \in V \cup E$,
\[
\mathrm{type}(x) \;\le_{\mathbf{Type}}\; \mathrm{type}'(f(x)),
\]
where $\le_{\mathbf{Type}}$ is the refinement relation in $\mathbf{Type}$.

\item \textbf{Tag consistency:}  
Tag refinements satisfy
\[
\mathrm{tag}(x) \;\le_{\mathbf{Tag}}\; \mathrm{tag}'(f(x)),
\]
and the projections match the type constraints.
\end{enumerate}
\end{definition}

\begin{example}
A vertex of type \texttt{Matter} with tag \texttt{Electron} may map to a vertex tagged \texttt{SubatomicParticle}, provided:
\[
\texttt{Electron} \le_{\mathbf{Tag}} \texttt{SubatomicParticle}
\quad\text{and}\quad
\texttt{Matter} \le_{\mathbf{Type}} \texttt{Matter}.
\]
\end{example}

Typed morphisms guarantee semantic safety under transformation.

% ============================================================
\section{Functorial Type Assignments}
% ============================================================

Type assignments extend functorially across Polyxan morphisms:
\[
\mathrm{type} : \mathbf{Polyxan} \to \mathbf{TypeCat},
\]
where $\mathbf{Polyxan}$ is the category of Polyxan hypergraphs with typed morphisms.

Functoriality implies:

\begin{itemize}
\item type preservation under composition:
\[
f, g \text{ typed} \;\Longrightarrow\; g \circ f \text{ typed},
\]
\item identity morphisms are trivially typed,
\item type-based constraints propagate through every rewrite and transformation.
\end{itemize}

This categorical foundation is essential for correctness of Polycompiler pipelines.

% ============================================================
\section{Typed Rewriting Rules}
% ============================================================

Typed rewriting rules constrain allowable transformations.

\begin{definition}[Typed Rewrite Rule]
A typed rewrite rule is a span of typed hypergraphs:
\[
L \xleftarrow{\ell} K \xrightarrow{r} R,
\]
where $L$ is the left-hand pattern, $R$ is the replacement, and $K$ is the interface.  
A rewrite may be applied to $\mathcal{G}$ when a typed morphism
\[
m: L \to \mathcal{G}
\]
exists that respects all structure, type, and tag constraints.
\end{definition}

Typed rewrites guarantee:
\begin{itemize}
\item \emph{type preservation} under substitution,
\item \emph{tag-order consistency} under refinement,
\item \emph{semantic stability} under Polyxan evolution,
\item safe rewrites in Polycompiler passes.
\end{itemize}

% ============================================================
\section{Applications to Polycompilation and Semantic Verification}
% ============================================================

Types and typed morphisms are central to:

\begin{itemize}
\item Polycompiler extraction/hoisting (Chapter~\ref{ch:polycompiler-algorithms}),
\item semantic invariants under deterministic or stochastic rewriting,
\item type-based constraint propagation across recursive hypergraphs,
\item coinductive proofs of infinite Polyxan processes (Chapter~\ref{ch:theorems}),
\item improved search, indexing, and retrieval of semantic atoms.
\end{itemize}

Typed rewriting provides the correctness backbone of the entire Polycompiler architecture.

% ============================================================
\section{Summary}
% ============================================================

This chapter established:

\begin{itemize}
\item a formal typing system for vertices and hyperedges,
\item tagging as a refinement structure with type projection,
\item typed morphisms preserving structure, type, and tag assignments,
\item a functorial perspective ensuring compositional invariants,
\item typed rewrite rules for safe, semantically consistent transformations,
\item and applications to Polycompiler correctness and verification.
\end{itemize}

This semantic type system forms the substrate for spans, co-spans, rewrite calculi, and self-rewriting atoms in forthcoming chapters.
